%% kapa: exam=1 round=28 year=2017 period=2 form=A subject=appraisal_law subject_ko=감정평가관계법규 seq=18 number=18
%% source: Q-Net 감정평가사 기출문제 (cst00309) artlSeq=5199603 file=extracted/1008/5199603/2017년 제28회 감정평가사 1차시험 문제지/2교시 A형.pdf
%% conversion: kordoc 4.24.51 -> tools/gen1.py (auto); review=proofread-sonnet; flags=none
\begin{question}{18}
부동산 가격공시에 관한 법령상 표준주택가격의 공시 등에 관한 설명으로 옳지 않은 것은?
\begin{choices}
  \item 국토교통부장관은 표준주택가격을 조사ㆍ산정하고자 할 때에는 한국감정원에 의뢰한다.
  \item 표준주택가격은 국토교통부장관이 중앙부동산가격공시위원회의 심의를 거쳐 공시하여야 한다.
  \item 표준주택의 대지면적 및 형상은 표준주택가격의 공시에 포함되어야 한다.
  \item 국토교통부장관이 표준주택가격을 조사ㆍ산정하는 경우에는 인근 유사 단독주택의 거래가격ㆍ임대료 등을 종합적으로 참작하여야 한다.
  \item 국토교통부장관은 개별주택가격의 산정을 위하여 필요하다고 인정하는 경우에는 주택가격비준표를 작성하여 시ㆍ도지사 또는 대도시 시장에게 제공하여야 한다.
\end{choices}
\end{question}
\answer{5}
